\section{通过训练获取推理能力}

\subsection{知识蒸馏：Orca}

提示方法只能用于大模型。如果我们想让\textbf{小模型}也具备推理能力，一个自然的思路是\textbf{知识蒸馏}（Distillation）：用大模型生成的推理解释来微调小模型。

Orca 是这一思路的典型代表。其流程分为三步：

\begin{enumerate}
    \item 从 FLAN V2 数据集收集大量指令（instructions）和问题
    \item 用 GPT-4 / ChatGPT 对这些指令生成\textbf{详细的解释性回答}（通过系统消息引导模型``请逐步解释你的推理过程''）
    \item 使用这些解释来微调 LLaMA-13B
\end{enumerate}

\begin{figure}[H]
\centering
\includegraphics[width=0.95\textwidth]{slides-images/slide-020.jpg}
\caption{Orca 的训练数据构建：从 FLAN V2 收集指令，共约500万样本\protect\footnotemark}
\end{figure}
\footnotetext{Slides 第21页。Source: Mukherjee et al. 2023。}

\subsubsection{评估：BigBench-Hard}

Orca 在 BigBench-Hard 上进行了评估。BigBench-Hard 包含23个聚焦多步推理的子任务，涵盖布尔表达式求值、日期理解、几何形状识别（根据 SVG 路径判断形状）等。

\begin{figure}[H]
\centering
\includegraphics[width=0.95\textwidth]{slides-images/slide-025.jpg}
\caption{BigBench-Hard 的日期理解任务示例：模型需要通过多步推理来确定``一年前的今天''的日期\protect\footnotemark}
\end{figure}
\footnotetext{Slides 第26页。Source: Suzgun et al. 2022。}

\begin{figure}[H]
\centering
\includegraphics[width=0.95\textwidth]{slides-images/slide-028.jpg}
\caption{BigBench-Hard 上 Orca 的表现：在多个子任务上优于 ChatGPT，也显著超过了同为 LLaMA-13B 的 WizardLM\protect\footnotemark}
\end{figure}
\footnotetext{Slides 第29页。}

\begin{importantbox}{蒸馏的有效性}
Orca 的成功说明：通过在大模型生成的\textbf{详细推理解释}上微调小模型，可以有效迁移推理能力。相比仅在（问题，答案）对上微调，加入中间推理步骤的训练数据能带来显著的性能提升。
\end{importantbox}

\subsection{自我训练：ReST$^{EM}$}

一个自然的问题是：为什么要用大模型的推理来训练小模型，能否让模型\textbf{在自己的推理输出上进行自我改进}？

ReST$^{EM}$（Reinforced Self-Training）正是基于这一思路，交替执行两个阶段：

\begin{enumerate}
    \item \textbf{E-Step（生成）}：给定推理问题，让模型\textbf{采样多个推理方案}，根据答案正确性进行\textbf{过滤}——仅保留得出正确答案的推理链
    \item \textbf{M-Step（更新）}：在过滤后的正确推理链上\textbf{微调模型}
\end{enumerate}

这两个步骤可以\textbf{迭代进行}：更新后的模型能生成更好的推理链，从而进一步改进模型。

\begin{figure}[H]
\centering
\includegraphics[width=0.95\textwidth]{slides-images/slide-030.jpg}
\caption{ReST$^{EM}$ 在 Hendrycks MATH 和 GSM8K 上的表现：随着迭代次数增加，性能逐步提升；橙色点为 PaLM-2-L，蓝色点为 PaLM-2-S，虚线为人类提供推理链的有监督微调基线\protect\footnotemark}
\end{figure}
\footnotetext{Slides 第31页（标注为29）。Source: Singh et al. 2024。}

\begin{importantbox}{自我训练可以超越人类标注}
ReST$^{EM}$ 的一个重要发现：经过多轮迭代自我训练，模型在自己生成的推理链上微调的效果可以\textbf{超越}在人类编写的推理链上微调的效果。这意味着模型自己生成的推理方式可能比人类提供的更适合模型自身的``认知模式''。
\end{importantbox}

\begin{warningbox}{自我训练的收敛问题}
在 GSM8K 上，ReST$^{EM}$ 的性能在2--3轮迭代后开始\textbf{下降}，表明自我训练并非总是单调改进的。可能的原因包括：（1）模型陷入自身偏见的``回声室''；（2）过滤器不够精确，部分错误推理链被保留；（3）数据多样性随迭代逐渐减少。
\end{warningbox}

\subsection{本章小结}

通过训练来获取推理能力有两条路径：（1）蒸馏——用大模型的推理解释微调小模型（如 Orca）；（2）自我训练——让模型在自己的正确推理上迭代微调（如 ReST$^{EM}$）。两种方法都表明，\textbf{推理链}（而非仅仅是正确答案）是训练推理能力的关键信号。

%% ========== Part 4: 推理的本质 ==========

